\documentclass[10pt]{article}
\usepackage[margin=1in]{geometry}
\usepackage{amsmath,amssymb}
\usepackage{hyperref}
\usepackage{booktabs}
\usepackage{listings}
\usepackage{xcolor}
\usepackage{microtype}
\usepackage{authblk}

\lstset{
  basicstyle=\ttfamily\footnotesize,
  breaklines=true,
  frame=single,
  backgroundcolor=\color{gray!8},
  columns=fullflexible
}

\title{\textbf{Equal Detector Error Models Do Not Make a QEC Compiler Pass Correct}}
\author{Cleiton Augusto Correa Bezerra}
\affil{Independent Researcher\\
\texttt{augusto.cleiton@gmail.com}}
\date{November 2026}

\begin{document}
\maketitle

\begin{abstract}
Comparing detector error models (DEMs) before and after a pass has been
proposed as a way to verify passes of a quantum error correction compiler. We
show a pass that inserts one deterministic Pauli and leaves the DEM identical
while changing what the circuit computes. On rotated surface codes of distance 3
and 5 this mutation goes unreported by the DEM comparison in all 200 trials,
and by self-referential detector sampling in all 200. A check on the circuit's
stabiliser flows reports all 200. The blindness also runs the other way: in a
repetition code, dropping a parity-check CX leaves the noiseless circuit
unchanged, so the flows accept it, while the DEM reports it in 100 of 100
trials. Neither check replaces the other. We package both as
\texttt{cleitonforge.qec} and run it over the circuit passes of Deltakit's two
compilers. Over 480 before-and-after pairs it finds no pass that changes a
circuit. It did find one defect, in our own DEM comparison.
\end{abstract}

\section{The claim}

A pass over an error-correction circuit should preserve two things: what the
circuit computes, and how its errors show up in the detectors. A recent preprint
gives a quasilinear decision procedure for DEM equivalence and proposes it for
this job, saying that such checks ``can be used to verify compiler
passes''~\cite{rennela2026}. It asks for nothing to be checked besides the DEM.
It notes that two equivalent circuits may produce different DEMs. It does not
consider the opposite case: two circuits that differ but produce the same DEM.

That case exists. Our earlier paper explains why: detectors and DEMs are defined
relative to the circuit's own noiseless execution~\cite{bezerra2026}. A
deterministic Pauli moves the measurement outcomes and that reference
together, so no error mechanism changes. The smallest instance:

\begin{lstlisting}[columns=fixed]
R 0                          R 0
X_ERROR(0.01) 0              X_ERROR(0.01) 0
M 0                          M 0
                             X 0
M 0                          M 0
DETECTOR rec[-1] rec[-2]     DETECTOR rec[-1] rec[-2]
\end{lstlisting}

On the right, the detector fires on every shot. Its DEM is identical to the one
on the left. So any procedure that compares only DEMs accepts this pass, however
exact it is. This is not an error in the decision procedure. It is a limit of
what a DEM records.

\section{The check}

\texttt{cleitonforge.qec} compares two circuits through their stabiliser flows,
in both directions, using Stim~\cite{gidney2021}: every flow generator of each
circuit must be a flow of the other. Noise is removed first, so the check is
about what the circuit computes. By default a reset is appended on every qubit,
since a memory experiment discards the final state. Without it, removing a
basis change left after the last measurement would be reported as a change,
which it is not.

Next to that it compares the DEMs as a decoder sees them. It checks detector and
observable counts, and the error mechanisms keyed by the symptoms they flip.
Mechanisms with the same symptoms are merged. Coordinates and line order are
ignored. A verdict carries both results and names the two disagreements:
\texttt{dem\_blind} (DEM equal, circuit changed) and \texttt{faults\_changed}
(circuit equal, DEM changed). The second is not always a bug, since some passes
change the DEM on purpose. It is always worth a look.

\begin{lstlisting}[language=Python]
from cleitonforge import qec   # pip install "cleitonforge[qec]"
v = qec.check_pass(before, after)  # stim.Circuit or deltakit Circuit
v.equivalent, v.dem_equal, v.dem_blind, v.faults_changed, v.witness
\end{lstlisting}

\section{How much each check sees}

We applied four mutations to Stim's generated circuits, 50 times each per
circuit: rotated surface code memory in $Z$ and $X$, and repetition code memory,
each at distance 3 and 5 with as many rounds. Noise was $10^{-3}$ after Clifford
gates, after resets and before measurements. A mutant counts as changed when
the flows say so. Because the flows are also one of the checks being measured,
we confirmed every verdict independently, without flows: the mutant's noiseless
shots, decoded with the original's detectors and observable, fire a detector or
flip the observable exactly when the flows report a change. The two agreed in
1200 of 1200 mutants.

\begin{table}[h]
\centering
\begin{tabular}{llrrrr}
\toprule
code & mutation & changed & DEM & detectors & flows \\
\midrule
surface & drop a two-qubit gate & 200 & 200 & 200 & 200 \\
surface & drop the control of a CX & 200 & 200 & 199 & 200 \\
surface & wrong measurement basis & 200 & 200 & 198 & 200 \\
surface & \textbf{X before a measurement} & 200 & \textbf{0} & \textbf{0} & 200 \\
repetition & drop the control of a CX & 100 & 100 & 0 & 100 \\
repetition & wrong measurement basis & 100 & 100 & 87 & 100 \\
repetition & \textbf{X before a measurement} & 100 & \textbf{0} & \textbf{0} & 100 \\
\bottomrule
\end{tabular}
\caption{Changed mutants each check reports. ``Detectors'' samples each circuit
against its own reference, as a detector sampler does. The flows column is
100\% by construction; the independent decode is what supports it.}
\end{table}

The fourth mutation, dropping a two-qubit gate, never changed a repetition code:
0 of 100. Every qubit there starts in $|0\rangle$ and a CX between such qubits
does nothing to the noiseless state. The independent decode agrees. But the DEM
differed in all 100 of those mutants, because a reset error the dropped CX used
to carry into an ancilla now reaches fewer detectors. This is the
\texttt{faults\_changed} case. The circuit computes the same thing, but its
errors show up differently, and a decoder built for the original would be
wrong about them. A flows check alone accepts it.

So the two checks fail on disjoint mutations. The DEM misses deterministic
Paulis; the flows miss changes that only matter under noise. A pass is
validated by both or by neither.

\section{Deltakit's passes}

We ran the check over the circuit passes of Riverlane's
Deltakit~\cite{deltakit}, in both of its compilers.

\paragraph{deltakit-explorer.} We compiled rotated and unrotated planar code
memories, $d=3$, in $Z$ and $X$, to three native gate sets. Then we checked
\texttt{compile\_circuit\_to\_native\_gates}, \texttt{remove\_identities},
\texttt{merge\_layers} (with and without breaking repeat blocks) and the new
\texttt{remove\_dead\_unitaries}. All 72 pairs are equivalent.
\texttt{remove\_dead\_unitaries} removes 20 gates over the 12 compiled circuits.

\paragraph{deltakit-compile.} Each pass ran between the Stim import and export
pipelines. The round trip alone, with no pass, is equivalent with an equal DEM
on every input. We used two sets of inputs: Stim's generated surface and
repetition codes at $d=3$ (24 passes, 72 pairs), and the 12 explorer
compilations with and without added noise (24 inputs, 14 passes, 336 pairs).
All 408 are equivalent, with equal DEMs. Three passes could not be
run on these inputs. One belongs at the Stim level and needs
\texttt{serialise-circuit} after it; at that level it was equivalent on the circuits we tried. One needs
arguments we did not supply. One declares that it does not support loops yet.

\paragraph{The defect we found.} The first run, with no pass at all, reported
the DEM as changed on every circuit. The import and export rewrite detector
coordinates, and our DEM comparison was textual. That was our defect, not
Deltakit's: a decoder ignores coordinates. It is fixed in 0.6.1, and it is the
reason the comparison now works on mechanisms.

\section{Limits}

Distances are 3 and 5, and all circuits are memory experiments; there is no
lattice surgery or adaptive circuit here. The four mutations are a small, fixed
set. The flows check needs Clifford circuits, which Stim requires anyway.
Finding no problem in Deltakit's passes on these inputs does not show that the
passes are correct in general.

\section{Reproducibility}

The tool is \texttt{cleitonforge} on PyPI, from 0.6.2. The table comes from
\texttt{research/qec\_pass\_validation.py} in the repository~\cite{cleitonforge},
with Stim 1.16 and seeded mutations. Its output is committed next to it.

\begin{thebibliography}{9}
\bibitem{rennela2026} M.~Rennela, \emph{Quasilinear Equivalence Checking for
Detector Error Models}, arXiv:2606.14677 (2026).
\bibitem{bezerra2026} C.~A.~C.~Bezerra, \emph{Detector-Based Oracles Are Empty
Where Quantum Error Correction Compilers Must Be Validated} (2026),
\href{https://doi.org/10.5281/zenodo.22802524}{doi:10.5281/zenodo.22802524}.
\bibitem{gidney2021} C.~Gidney, \emph{Stim: a fast stabilizer circuit
simulator}, Quantum 5, 497 (2021).
\bibitem{deltakit} Riverlane, \emph{Deltakit},
\url{https://github.com/Deltakit}.
\bibitem{cleitonforge} C.~A.~C.~Bezerra, \emph{CleitonForge},
\url{https://github.com/cleitonaugusto/CleitonForge}.
\end{thebibliography}

\end{document}
